% Pursuing Stacks (cotes 134-2 to 134-8): archive cote 134-7, site page 97 (PDF page 98), typescript page 539.
% typescript. Draft compared: ps6.tex:4885-4965.
% AI-transcribed by Claude Opus 5.5 (claude-opus-5-5, adaptive thinking, effort high), Message Batches API, batch ps4 (msgbatch_012Aoc1NuskATPsyBeKSYK7J), 2026-11-12; not yet checked by eye.
\scanpage{134-7/97}
\mspage{539}

Reverting to a general notion of ``bundles'' U(k),  we define now a
linearization functor
\[ (11) \qquad L:\ U(k) \longrightarrow Uab(k) \]
to be a functor left adjoint to the forgetful functor (7),  i.e. giving
rise to an adjunction isomorphism
\[ (12) \qquad \xout{Hom(X} \ \mathrm{Hom}_{U(k)}(X, K(L)) \simeq \mathrm{Hom}_{Uab(k)}(L(X), L) \quad , \]
where X is in U(k),  L in Uab(k).  Passing over to semisimplicial objec\unsure{\hand{ts}}
and the corresponding derived categories,  the functors K and L in
(7) and (11) should give rise to the functors $\overset{\hand{$\sim$}}{K}$ and LH.  of section
124.

When we take  $U(k) = U_W(k)$,  there is a drawback though,  because
(as we saw in section 115) the functor K \uline{does not} admit a left adjoint,
only a \uline{proadjoint},  associating to an object X a\xout{n} \uline{proobject} L(X) of
Uab(k).  As pointed out in section 124 (p. 522),  if we want a nice
pair (K,L) of adjoint functors,  this forces us to work with prounipo\xout{t}-%
tent bundles instead of just unipotent ones,  thus getting out of the
haven of sheaves over k and into the somewhat dubious sea of prosheaves
and semisimplicial prosheaves,  which have not been provided for in
Illusie's foundational ponderings !  If we do stick to the W-approach,\margin{we \unsure{can} however work with the \unsure{$\varprojlim$} of \unsure{corr.} sheaves ---}
this promises \xout{with} \inserted[typed]{us} a fair amount of extra sweat,  putting in proobjects
everywhere,  not tooenticing a prospect is it ?

It would seem that we are better off with the V-approach,  in which
case the \xout{objects} bundles (more accurately, V-bundles) we are working \xout{wi}
with are actual schemes,  indeed affine schemes over k,  as we get
\[ (13) \qquad V(M) \simeq \mathrm{Spec}(\mathrm{Sym}_k(M)) \quad . \]
Now \unsure{let|lat} more generally X be any scheme,  and let's look at maps from
X into any vector bundle V(M),  we get
\[ \hand{(*)} \qquad \mathrm{Hom}(X,V(M)) \overset{\hand{def}}{=} \mathrm{Hom}_{\underline{O}_{\hand{X}}}(p^*(\widetilde{M}\,), \underline{O}_X) \overset{\hand{adjunction}}{\simeq} \mathrm{Hom}_{\underline{O}_{\over{S}{S}}}(\widetilde{M}\,, p_*(\underline{O}_X)\,)\ , \]
where
\[ p:\ X \longrightarrow S = \mathrm{Spec}\ k \]
is the structural map of X,\inserted{and $\widetilde{M}$ the restriction of W(M) to the usual Zariski site of S = Spec(k).} It is well known that under a rather mild
restriction on X (namely X \deleted{being} quasi-compact and quasi-separated),
always satisfied when X is affine,  $p_*$ takes quasi-coherent sheaves on
X (for the usual small Zariski site) into quasi-coherent sheaves on
S = Spe\unsure{c}\xout{x}k ;  thus \hand{if}

\margin{97}
